\documentclass[border=0pt]{standalone}
% The shared preamble of every docs diagram (docs/diagrams/<name>.tex).
% dvisvgm's driver first, so patterns and opacity survive into the SVG.
\def\pgfsysdriver{pgfsys-dvisvgm.def}
\usepackage[siunitx,RPvoltages]{circuitikz}
\usepackage{xcolor}
% Latin Modern: outline fonts for every glyph a label uses, µ included.
\usepackage[T1]{fontenc}
\usepackage{lmodern}
\usetikzlibrary{arrows.meta,backgrounds,calc,patterns,positioning}
\tikzset{
  % A white card, so diagrams read in mdBook's dark themes too.
  every picture/.append style={show background rectangle,
    background rectangle/.style={fill=white,inner sep=4pt}},
  % Layout geometry: layers told apart by fill and hatching, never colour alone.
  layer a/.style={draw=black,line width=0.6pt,fill=gray!20},
  layer b/.style={draw=black,line width=0.6pt,fill=gray!45},
  layer c/.style={draw=black,line width=0.6pt,fill=white,postaction={pattern=north east lines}},
  layer d/.style={draw=black,line width=0.6pt,fill=white,postaction={pattern=crosshatch dots}},
  marker/.style={draw=black,dashed,line width=0.6pt},
  % A measured distance and its label.
  dim/.style={{Latex[length=2mm]}-{Latex[length=2mm]},line width=0.5pt},
  dim label/.style={font=\small,fill=white,inner sep=1pt},
  % Where a rule fails: thick, dashed and labelled, not only red.
  violation/.style={draw=red!75!black,line width=1.4pt,dash pattern=on 3pt off 2pt},
  violation label/.style={font=\small\bfseries,text=red!75!black},
  caption/.style={font=\small},
}

\begin{document}
\begin{circuitikz}[x=1.4cm, y=1.4cm]
  % cross-section: a 1 um met1 wire 1.376 um above the substrate; its heat
  % spreads down through the dielectric over W + phi t
  \draw[layer b] (-3,-0.5) rectangle (4,0);
  \node[caption] at (3.3,-0.25) {substrate};
  \draw[layer a] (0,1.376) rectangle (1,1.736);
  \node[caption, anchor=west] at (1.1,1.556) {met1, $W = \SI{1}{\micro\metre}$, carrying \SI{10}{\milli\ampere}};
  \draw[dashed] (0,1.376) -- (-0.605,0);
  \draw[dashed] (1,1.376) -- (1.605,0);
  \draw[dim] (-0.4,0) -- (-0.4,1.376);
  \node[dim label, anchor=east] at (-0.5,0.69) {$t = \SI{1.376}{\micro\metre}$};
  \draw[dim] (-0.605,-0.75) -- (1.605,-0.75);
  \node[dim label, below] at (0.5,-0.8) {$W + \varphi t = \SI{2.21}{\micro\metre}$};
  \node[caption, align=left, anchor=north west] at (-3,-1.35) {%
    $\Delta T = I^2 R_\mathrm{s}\, t / \bigl(K W (W + \varphi t)\bigr) = \SI{5.56}{\kelvin}$,
    so the wire runs at \SI{110.6}{\celsius}, not \SI{105}{\celsius}};
  \node[violation label, align=left, anchor=north west] at (-3,-1.85) {%
    allowed at \SI{110.6}{\celsius}: \SI{9397}{\micro\ampere} $<$ \SI{10000}{\micro\ampere} carried
    (\SI{10978}{\micro\ampere} at \SI{105}{\celsius} would pass)};
\end{circuitikz}
\end{document}
